% Supplemental Material for "Exact exponential heat rectification from the Liouvillian skin effect"
\documentclass[aps,prl,superscriptaddress,nofootinbib,onecolumn,notitlepage,11pt]{revtex4-2}
\usepackage{amsmath,amssymb,bm}
\usepackage{graphicx}
\graphicspath{{figures/}}
\usepackage{booktabs,multirow}
\usepackage[colorlinks=true,linkcolor=blue,citecolor=blue,urlcolor=blue]{hyperref}
\newcommand{\R}{\mathcal{R}}
\newcommand{\dQ}{\delta Q}
\newcommand{\heff}{h_{\rm eff}}
\setcounter{secnumdepth}{2}
\renewcommand{\thesection}{S\arabic{section}}
\renewcommand{\thetable}{S\arabic{table}}
\renewcommand{\theequation}{S\arabic{equation}}

\begin{document}

\title{Supplemental Material for ``Exact exponential heat rectification from the Liouvillian skin effect''}
\author{Hikaru Wakaura}
\affiliation{QIRI (Quantum Integrated Research Institute Inc.), Tokyo 107-0061, Japan}
\author{Taiki Tanimae}
\affiliation{QIRI (Quantum Integrated Research Institute Inc.), Tokyo 107-0061, Japan}
\maketitle

\section{Gaussian solution and the $\Delta T$ response}
\label{sec:gauss}

The Lindbladian of the main text, $\dot\rho=-i[H,\rho]+\sum_a\mathcal{D}[L_a]\rho$, is quadratic for $V=0$: $H=\sum_{mn}h_{mn}c_m^\dagger c_n$ and every jump operator is linear, $L_a=\sum_m l_{a,m}c_m$ (loss) or $L_a=\sum_m g_{a,m}c_m^\dagger$ (gain). The leads are included explicitly as sets of modes with their own loss and gain jumps whose rates realise a thermal occupation. The two-point function $C_{mn}=\langle c_m^\dagger c_n\rangle$ then obeys a closed equation,
\begin{equation}
\dot C=-(XC+CX^\dagger)+M_g^{T},\qquad X=\tfrac12(M_l+M_g^{T})-ih^{T},\qquad M_l=\sum_a l_al_a^\dagger,\quad M_g=\sum_a g_ag_a^\dagger,
\end{equation}
whose steady state is the solution of the Lyapunov equation $XC+CX^\dagger=M_g^{T}$. We solve it with a Bartels--Stewart solver (\texttt{scipy.linalg.solve\_continuous\_lyapunov}) and verified it against the exact steady state of the full many-body Lindbladian on four modes to $10^{-8}$. The energy current from reservoir $\alpha$ into the system is obtained from the jump operators labelled $\alpha$ and from $h$ and $C$; the single-particle amplitudes propagate with $\heff=h-\tfrac{i}{2}(M_l^{T}+M_g)$, which for pure bond loss is Eq.~(2) of the main text.

\paragraph{Leads.}
Two implementations are used. (i) \emph{Wide-band leads}: the end sites acquire a self-energy $-i\Gamma/2$ and the heat currents are given by the Landauer formula with $T_{RL}(\varepsilon)=\Gamma^2|G_{L1}(\varepsilon)|^2$, $T_{LR}(\varepsilon)=\Gamma^2|G_{1L}(\varepsilon)|^2$, $G=(\varepsilon-\heff)^{-1}$; we use $\Gamma=J$ and integrate over $|\varepsilon|\le4J$. (ii) \emph{Mesoscopic leads}: $K$ modes equally spaced in $[-W,W]$, each coupled to its end site with $v=\sqrt{\Gamma\Delta\varepsilon/2\pi}$ and thermalised by loss and gain rates $\gamma(1-f)$ and $\gamma f$, with $f$ the Fermi function at $(T_\alpha,\mu)$ and $\gamma=\Delta\varepsilon$ chosen to make the lead spectrum smooth. We use $K=10$, $W=4J$ for the matrix-product-state (MPS) runs and $K=10$--$80$ for checks. The $\Delta T$ responses of the two implementations agree within $5$--$15\%$, the difference being a bandwidth effect; the ratio $\dQ_f/\dQ_r$ obeys Eq.~(3) exactly in both, since only the end self-energies differ and they are diagonal.

\paragraph{Why total currents are the wrong observable.}
With bulk loss, the total heat current absorbed by the cold lead contains a contribution that does not depend on the temperature of the other lead: both leads lose energy to the loss channel. If rectification is defined from total currents, this background dominates and the apparent $\R$ decays as $1/L$ for any $\kappa$, although the device is perfectly directional. The $\Delta T$ responses $\dQ_{f}=\dot E_R(T_h,T_c)-\dot E_R(T_c,T_c)$ and $\dQ_r=\dot E_L(T_c,T_h)-\dot E_L(T_c,T_c)$ (heat absorbed by the respective cold lead) remove the background exactly and are the quantities whose ratio obeys the exact law. Throughout, $T_h=J$, $T_c=0.5J$ and $\mu=0$ unless stated otherwise; $\hbar=k_B=1$.

\section{Proofs of the exact statements}
\label{sec:proofs}

\paragraph{Tridiagonal identity.}
For a tridiagonal matrix $A$ of dimension $L$ with sub-diagonal $a_{j+1,j}$ and super-diagonal $a_{j,j+1}$, the corner elements of the inverse are $(A^{-1})_{L1}=(-1)^{L+1}\prod_{j}a_{j+1,j}/\det A$ and $(A^{-1})_{1L}=(-1)^{L+1}\prod_j a_{j,j+1}/\det A$ (expand the cofactor of the corner element; the only non-vanishing term is the product of the off-diagonal chain). With $A=\varepsilon-\heff$ this gives $G_{L1}/G_{1L}=\prod_j(\heff)_{j+1,j}/(\heff)_{j,j+1}$. The product involves no diagonal element, so it is independent of $\varepsilon$, of the on-site energies, of the end self-energies and of the uniform loss. For the uniform chain all factors equal $t_-/t_+$ and $T_{RL}/T_{LR}=|t_-/t_+|^{2(L-1)}$ at every energy. The thermal averages $\dQ_{f,r}=\int\frac{d\varepsilon}{2\pi}\varepsilon T_{RL,LR}(\varepsilon)\Delta f(\varepsilon)$ share the same weight, hence Eq.~(3). The argument does not use the fermionic nature of the leads: for bosonic carriers with the same bond loss $\Delta f$ is a difference of Bose functions and the ratio is unchanged.

\paragraph{Bond disorder and the sum rule.}
With bond-dependent $\kappa_j$ and $\phi_j$ the factors are $t_{\mp,j}=-J-\tfrac{i\kappa_j}{2}e^{\mp i\phi_j}$ and the identity gives Eq.~(5) of the main text directly. We verified the sum rule to $10^{-14}$ on $10^2$ random samples, with on-site disorder present at the same time.

\paragraph{Same-chirality gain.}
Adding $G_j=\sqrt{\kappa g}\,(c_j^\dagger+e^{-i\phi}c_{j+1}^\dagger)$ produces $M_g$ with the same matrix elements as $M_l^{T}$, so $\heff$ is obtained from the pure-loss one by $\kappa\to\kappa(1+g)$: the law keeps its form with a shorter $\xi$. Consequently a local bath at infinite temperature and $\mu_b=0$, implemented as loss $\kappa/2$ plus gain $\kappa/2$ of the same chirality, has exactly the free-fermion ratio of a pure loss $\kappa$. This is the configuration used for the green symbols in Fig.~3 of the main text.

\paragraph{Mirror reflection equals transposition and the identity $\sigma_f(T_b=T_c)=\sigma_f(T_b=T_h)$.}
Let $P$ be the site reflection $j\to L+1-j$. For the uniform chiral chain $P\heff P$ exchanges $t_+$ and $t_-$, which is also what transposition does, so $P\heff P=\heff^{T}$ and $P G P=G^{T}$. The end reflectances therefore coincide, $G_{LL}=(PGP)_{11}=(G^{T})_{11}=G_{11}$, although $G_{L1}\neq G_{1L}$. In the multi-terminal Landauer form with all reservoirs thermal, the entropy production with every reservoir except $L$ at $T_c$ is $\sigma=(1/T_c-1/T_h)\int\frac{d\varepsilon}{2\pi}\varepsilon\,[f_h-f_c]\,[1-R_{LL}(\varepsilon)]$, and with every reservoir except $R$ at $T_h$ it is the same expression with $R_{RR}$ (the sum rules $\sum_\beta T_{\alpha\beta}=\sum_\beta T_{\beta\alpha}=1-R_{\alpha\alpha}$ are used). Equal reflectances therefore give $\sigma_f(T_b=T_c)=\sigma_f(T_b=T_h)$ for any $\phi$, $\kappa$ and $t_2$. On-site disorder breaks the identity (difference $3\times10^{-3}$ for $W_d=J$), and mirror-symmetric disorder restores it ($10^{-12}$); the clean difference is $10^{-12}$. Physically, the heat emitted by a hot lead into an otherwise cold device equals the heat absorbed by a cold lead from an otherwise hot device: nonreciprocity appears in transmission but not in reflection.

\section{Non-Bloch form of the exponent for next-nearest-neighbour hopping}
\label{sec:nonbloch}

With $t_2\neq0$ the bulk ansatz $\psi_j=\beta^j$ turns $\heff\psi=\varepsilon\psi$ into the quartic Eq.~(4) of the main text. Its roots satisfy $\beta_1\beta_2\beta_3\beta_4=1$ by Vieta's formula. For a long open chain the transfer matrix from site 1 to site $L$ is dominated by the two roots of intermediate modulus, which form the generalised Brillouin zone when $|\beta_2|=|\beta_3|$ and otherwise split it; the amplitude transported from left to right scales as $|\beta_3|^{L}$ and from right to left as $|\beta_2|^{-L}$, giving $\ln[T_{RL}/T_{LR}](\varepsilon)=2(L-1)\ln|\beta_2\beta_3|(\varepsilon)+O(1)$. For $t_2=0$ the equation is quadratic with $\beta_1\beta_2=t_-/t_+$, and the $O(1)$ term vanishes identically. Figure~2(a) of the main text compares the prediction with the exact transmissions at $\kappa=0.2J$; the agreement is equally good at $\kappa=0.1J$ and $0.3J$ for all $t_2\le0.5J$.

The periodic-boundary spectrum is $E(k)=-2J\cos k+2t_2\cos2k-i\kappa(1-\sin k)$ for $\phi=\pi/2$. Its imaginary part is non-positive and vanishes only at $k=\pi/2$, so the loop touches the real axis at one point and the winding number $w(\varepsilon)=\frac{1}{2\pi i}\oint d\ln[E(k)-\varepsilon]$ is zero for every real $\varepsilon$ and every $t_2$. The winding around the centre of the loop, $-i\kappa$, is $-1$ and signals the skin effect, but it is a single integer and cannot encode the energy dependence, sign change and $t_2$ dependence of the transmission asymmetry; the quantity that does is $\ln|\beta_2\beta_3|(\varepsilon)$.

\section{Superfermion matrix-product-state method and its precision}
\label{sec:mps}

\paragraph{Vectorisation.}
We use the superfermion representation: $|\rho\rangle\!\rangle=\rho|I\rangle\!\rangle$ with $|I\rangle\!\rangle=\prod_j(1+c_j^\dagger\tilde c_j^\dagger)|0\rangle$, where $\tilde c_j$ are auxiliary (bra) fermions. For a parity-even $\rho$ the rules $c_j\rho\to c_j$, $\rho c_j^\dagger\to-\tilde c_j$, $\rho c_j\to\tilde c_j^\dagger$, $\rho c_i^\dagger c_j\to\tilde c_j^\dagger\tilde c_i$ map the Lindbladian onto a local operator on $2N$ fermionic modes; the bond dissipators $L_j\rho L_j^\dagger$ become bilinears on four modes and the interaction $Vn_jn_{j+1}$ appears as $V n_jn_{j+1}-V\tilde n_j\tilde n_{j+1}$. Each physical site is grouped with its auxiliary partner into one MPS site of local dimension four, and the conserved charge $N-\tilde N=0$ is exploited. Expectation values are $\langle\!\langle I|O|\rho\rangle\!\rangle/\langle\!\langle I|\rho\rangle\!\rangle$. The lead jump operators commute with the interaction, so the lead energy currents are two-point functions and the Gaussian current formula applies unchanged.

\paragraph{Layout and propagation.}
The $L$ chain sites and the $2K$ lead modes are arranged as [site 1][left lead][sites $2\dots L$][right lead], which keeps all long-range couplings anchored at a chain end and the MPO bond dimension of the Lindbladian at six. The state is propagated with the second-order $W^{\rm II}$ MPO exponential of TeNPy in steps $dt$, with SVD compression at bond dimension $\chi$ and singular-value cutoff $10^{-12}$, starting from the product state with the exact Gaussian occupations, until the change of both lead currents over a chunk of $2/J$ is below $10^{-5}$ or $t=60/J$ is reached. The mapping was verified exactly: on four modes with $V=0.7J$ the null vector of the MPO Lindbladian coincides with the many-body steady state to $10^{-15}$, and for $V=0$, $L=6$, $K=4$ the propagated steady state agrees with the Lyapunov solution to $2\times10^{-3}$ in the currents.

\paragraph{Precision.}
The reverse response $\dQ_r$ is exponentially small and is the limiting quantity. The truncation error of the absolute currents accumulates with the number of time steps, so it is \emph{not} reduced by a smaller $dt$; it is reduced only by $\chi$. Table~\ref{tab:chi} lists the free-fermion calibration at $\kappa=0.05J$, $L=16$. At $\chi=256$ and $dt=0.05/J$ the ratio is reproduced within $0.003$ and $\dQ_r$ within $0.03\%$; this setting is used for all production runs. At $\kappa=0.1J$ the residual bias of $\ln(\dQ_f/\dQ_r)$ at $\chi=256$ is $+0.045$ ($L=12$) and $-0.038$ ($L=16$), which we quote as the error bar; slopes at this $\kappa$ were extracted from the control-variate estimate $\ln r_{\rm cv}(V)=\ln r_{\rm MPS}(V)-\ln r_{\rm MPS}(0)+\ln r_{\rm exact}(0)$, which cancels the common truncation bias. The time-fluctuation of the last five chunks contributes less than $0.01$ to $\ln(\dQ_f/\dQ_r)$ except for the infinite-temperature-bath run at $V=2J$, where it is $0.05$.

\begin{table}[h]
\caption{Free-fermion calibration of the MPS steady state at $\kappa=0.05J$, $L=16$, $K=10$: logarithmic ratio and reverse response compared with the exact Gaussian values $\ln(\dQ_f/\dQ_r)=1.5003$ and $\dQ_r=4.918\times10^{-3}J^2/\hbar$.}
\label{tab:chi}
\begin{ruledtabular}
\begin{tabular}{lcccc}
 & $\chi=128$, $dt=0.1$ & $\chi=128$, $dt=0.05$ & $\chi=128$, $dt=0.025$ & $\chi=256$, $dt=0.05$ \\
\colrule
$\ln(\dQ_f/\dQ_r)$ & 1.347 & 1.412 & 1.462 & 1.5033 \\
$\dQ_r$ bias & $+16.8\%$ & $+7.7\%$ & $+2.6\%$ & $+0.03\%$ \\
\end{tabular}
\end{ruledtabular}
\end{table}

\paragraph{Resolvable regime.}
Because the MPS floor of the absolute currents is $\sim10^{-3}J^2/\hbar$ at $\chi=128$ and about two orders of magnitude lower at $\chi=256$, the reverse response can be resolved only when $\dQ_r/\dQ_f=e^{-2(L-1)/\xi}$ is not much smaller than $10^{-2}$, i.e.\ $2(L-1)\kappa/J\lesssim4$. For stronger skin effect (e.g.\ $\kappa=0.2J$, $L=20$--$40$) the simulations yield only lower bounds, $\R\ge0.85$--$0.98$ for all $V\le2J$, together with the forward throughput, which decreases monotonically with $V$ ($1.10\to0.63\times10^{-2}J^2/\hbar$ at $L=20$ and $0.70\to0.33\times10^{-2}$ at $L=40$ for $V=0\to2J$).

\paragraph{Interacting results.}
Table~\ref{tab:int} collects the $\chi=256$ values used in Fig.~3. The $\chi=128$ series at $\kappa=0.03J$ and $0.05J$ over $L=12,16,20$ give slopes $2/\xi_{\rm eff}$ of $0.055,0.066,0.046,0.043\pm0.013$ ($\kappa=0.03J$; exact $0.060$) and $0.096,0.101,0.125,0.094\pm0.019$ ($\kappa=0.05J$; exact $0.100$) for $V=0,0.5,1,2J$. A direct linear-response evaluation of $\dQ_{f,r}$ by integrating the perturbed Lindbladian dynamics was also implemented; it is exact for $N\le14$ modes but, at equal $\chi$, less accurate than the subtraction method for $N\ge32$, and was not used.

\begin{table}[h]
\caption{$\chi=256$ steady-state results ($dt=0.05/J$, $K=10$, $\phi=\pi/2$). ``bath'' denotes the local infinite-temperature bath (loss $\kappa/2$ + gain $\kappa/2$, total rate $0.1J$).}
\label{tab:int}
\begin{ruledtabular}
\begin{tabular}{lcccccc}
$\kappa/J$ & $L$ & exact $V=0$ & MPS $V=0$ & MPS $V=2J$ & $\Delta_V$ & slope $V=2J$ / exact \\
\colrule
0.05 & 16 & 1.5003 & 1.5033 & 1.083 & $-0.42$ & \multirow{2}{*}{$0.0975/0.1000$} \\
0.05 & 20 & 1.9004 & --- & 1.473 & $-0.43$ & \\
0.1 & 12 & 2.2018 & 2.247 & 1.704 & $-0.54$ & \multirow{2}{*}{$0.207/0.200$} \\
0.1 & 16 & 3.0025 & 2.964 & 2.448 & $-0.52$ & \\
0.1 (bath) & 12 & 2.2018 & 2.203 & 1.662 & $-0.54$ & --- \\
\end{tabular}
\end{ruledtabular}
\end{table}

\section{Finite-temperature bulk bath}
\label{sec:bath}

Each bond $j$ is coupled to $K_b$ bath modes with energies in $[-W_b,W_b]$, thermalised at $(T_b,\mu_b)$ as in the leads, through the bond operator $A_j=c_j+e^{i\phi}c_{j+1}$ with amplitude $v_b=\sqrt{\kappa\Delta\varepsilon_b/2\pi}$, so that the damping rate of $A_j$ is $\kappa$ and the pure-loss model is recovered for $\mu_b\to-\infty$. All reservoirs now have a temperature and $\sigma=-\sum_\alpha(\dot E_\alpha-\mu_\alpha\dot N_\alpha)/T_\alpha$ is defined; we verified $\sigma\ge0$ in all configurations and $\sigma=0$ to $10^{-8}$ when all reservoirs share one temperature. We use $L=40$, $K_{\rm lead}=10$, $K_b=10$, $W_b=4J$.

Because the bath self-energy is energy dependent, the effective $t_\pm(\varepsilon)$ are not constant and the ratio falls slightly below the wide-band law: at $\kappa=0.2J$ we obtain $\ln(\dQ_f/\dQ_r)=14.23$, $14.48$, $14.68$ and $14.95$ for $(K_b,W_b/J)=(10,4)$, $(10,6)$, $(20,4)$ or $(20,6)$, and $(40,6)$, approaching the wide-band value $15.65$. The Lindblad loss is the exact wide-band limit. The ratio is identical for all $(T_b,\mu_b)$ at fixed discretisation (difference $<10^{-9}$).

With $T_b=T_c$ and $\mu_b=0$ the isothermal entropy production vanishes and $\sigma_f=0.042$, $0.044$, $0.047$, $0.049$, $0.052\,J/\hbar$ for $\kappa=0.05$, $0.1$, $0.2$, $0.3$, $0.5J$, while $\dQ_f=1.64$, $1.11$, $0.66$, $0.47$, $0.29\times10^{-2}J^2/\hbar$; the cost $\sigma_f/\dQ_f=2.6$, $4.0$, $7.0$, $10.4$, $18.1\,J^{-1}$ is the quantity in Fig.~4(a). The cost is carried by the heat flowing from the hot lead into the bath, $\dot Q_{{\rm bulk},f}=0.026\to0.049\,J^2/\hbar$ over the same range. For $\kappa=0.2J$ the $\mu_b$ sweep gives $\sigma_f=0.047$, $0.12$, $0.34$, $1.15$, $3.2$, $25.5\,J/\hbar$ and mean filling $0.50$, $0.44$, $0.38$, $0.25$, $0.12$, $0.10$ for $\mu_b=0$, $-0.5$, $-1$, $-2$, $-4$, $-30J$; the skin-like density pile-up at the edge appears only for $\mu_b<0$. For $T_b\neq T_c$ the isothermal state already produces entropy through bath--lead heat flow and $\sigma_f-\sigma_{\rm iso}$ can be negative; the cost indicator is therefore quoted at $T_b=T_c$.

\section{Additional disorder data}
\label{sec:disorder}

On-site disorder $\varepsilon_j\in[-W_d/2,W_d/2]$ (uniform), $L=40$, wide-band leads, $60$--$100$ samples per point. The rectification is sample independent to $10^{-14}$: $\R=0.9604$, $0.9992$, $1.0000$ for $\kappa=0.05$, $0.1$, $0.2J$ at every $W_d$. The median forward throughput at $\kappa=0.1J$ is $1.62$, $1.26$, $0.61$, $0.051$, $0.0014$, $0.00016\times10^{-2}J^2/\hbar$ for $W_d=0$, $0.5$, $1$, $2$, $3$, $4J$. The forward decay rate $-d\langle\ln\dQ_f\rangle/dL$ between $L=10$ and $80$ is listed in Table~\ref{tab:decay}; in the weak-disorder regime the loss-induced and Anderson decay rates are nearly additive, and no ``delocalisation by the skin effect'' appears in the forward direction, which is attenuated by loss and localisation alike. Bond disorder at $\kappa=0.1J$: a spread $\kappa_j=\kappa(1\pm0.8u)$, $u$ uniform in $[0,1]$, gives means $3.75$, $7.85$, $15.83$ (clean $3.80$, $7.81$, $15.81$) and standard deviations $0.43$, $0.55$, $0.88$ for $L=20$, $40$, $80$; a phase spread $\phi_j=\pi/2\pm1$ gives means rescaled by $\langle\sin\phi_j\rangle=\sin1\simeq0.84$ ($6.56$ against $7.81\times0.84=6.57$ for $L=40$) and standard deviations $0.13$--$0.25$; random chirality $\phi_j=\pm\pi/2$ gives means $-0.07$, $-0.04$, $0.16$ and standard deviations $0.91$, $1.34$, $1.77$. With $t_2=0.3J$ and $\kappa=0.2J$ the ratio becomes sample dependent, $\ln(\dQ_f/\dQ_r)=10.12\pm0.11$, $9.71\pm0.19$, $9.42\pm0.30$ for $W_d=0.5,1,2J$ (clean $10.31$).

\begin{table}[h]
\caption{Forward decay rate $-d\langle\ln\dQ_f\rangle/dL$ (per site, $L=10\to80$) for on-site disorder of width $W_d$.}
\label{tab:decay}
\begin{ruledtabular}
\begin{tabular}{lccc}
 & $W_d=0.5J$ & $W_d=J$ & $W_d=2J$ \\
\colrule
$\kappa=0$ & 0.006 & 0.020 & 0.055 \\
$\kappa=0.05J$ & 0.017 & 0.034 & 0.090 \\
$\kappa=0.2J$ & 0.027 & 0.044 & 0.107 \\
\end{tabular}
\end{ruledtabular}
\end{table}

\section{Experimental parameters and related work}
\label{sec:exp}

Table~\ref{tab:exp} summarises the parameter sets quoted in the main text. The requirement $\R>0.99$ corresponds to $(L-1)/\xi>2.65$ with $\xi\simeq J/\kappa$ at $\phi=\pi/2$. Related but distinct results are the diffusive skin effect in classical thermal metamaterials~\cite{Cao2021}, which concerns heat conduction rather than quantum transport between reservoirs; non-reciprocal routing of thermal noise in few-mode cascaded networks~\cite{Xuereb2018}, which does not scale with system size; the transport theory of non-Hermitian nonreciprocal systems~\cite{Ghaemi2021}, which treats reflectionless and transparent transport but not the thermal rectification ratio; and the recent photonic observation of skin dynamics in the Liouvillian regime~\cite{Yang2026}. For a review of boundary-driven open quantum systems see Ref.~\cite{Landi2022}.

\begin{table}[h]
\caption{Parameter sets for the two proposed platforms.}
\label{tab:exp}
\begin{ruledtabular}
\begin{tabular}{lcc}
 & cold atoms (momentum lattice) & superconducting circuit \\
\colrule
$J/2\pi$ & 0.5 kHz & 10--50 MHz \\
auxiliary damping $\Gamma_a/2\pi$ & optical loss & $\sim10$ MHz \\
coupling $g/2\pi$ & Raman & 1--3 MHz \\
$\kappa=4g^2/\Gamma_a$, $\kappa/2\pi$ & 25--100 Hz & 0.5--2 MHz \\
$\kappa/J$ & 0.05--0.2 & 0.05--0.2 \\
$L$ for $\R>0.99$ & 15 ($\kappa/J=0.2$) to 60 ($0.05$) & 15 to 60 \\
thermal leads & two-terminal reservoirs~\cite{Brantut2012,Krinner2017} & noise temperature 20--300 mK on end lines \\
carriers & fermionic atoms & photons (law unchanged, Sec.~\ref{sec:proofs}) \\
\end{tabular}
\end{ruledtabular}
\end{table}

\bibliographystyle{apsrev4-2}
\bibliography{refs}

\end{document}
